\subsection{导子的泛问题；交换情形}

现在假设 A 是一个交换 K-代数，E 是一个 A-模；可以把 E 看作一个 (A, A)-双模，其两个外合成律都与给定的 A-模运算相同。另一方面，\(A \otimes_{K} A\) 上的 (A, A)-双模结构与其 \((A \otimes_{K} A)\) -模结构相同，后者源自 \(A \otimes_{K} A\) 上的交换环结构，因为在这种情形下，对于 A 中的 x、y、u、v，

\[
x. (u \otimes v). y = (x u) \otimes (v y) = (x u) \otimes (y v) = (x \otimes y) (u \otimes v).
\]

因此在这种情形下，\(\mathfrak{J}\) 作为 \(m\) 的核是环 \(\mathbf{A} \otimes_{\mathbb{K}} \mathbf{A}\) 的理想，并且，由于 \(m: \mathbf{A} \otimes_{\mathbb{K}} \mathbf{A} \to \mathbf{A}\) 是满射的， \((\mathbf{A} \otimes_{\mathbb{K}} \mathbf{A}) / \mathfrak{J}\) 同构于 \(\mathbf{A}\) ；若还把 \(\mathbf{E}\) 视为一个 \((\mathbf{A} \otimes_{\mathbb{K}} \mathbf{A})\) -模，即经由 \(m\) （换言之，即 \((\mathbf{A} \otimes_{\mathbb{K}} \mathbf{A})\) -模 \(m_*(\mathbf{E})\) ），则 \((\mathbf{A}, \mathbf{A})\) -双模同态 \(\mathfrak{J} \to \mathbf{E}\) 便与 \((\mathbf{A} \otimes_{\mathbb{K}} \mathbf{A})\) -模同态 \(\mathfrak{J} \to \mathbf{E}\) 相等同 ( \(\S 4, no. 3\) )，换言之，存在一个典范的 \(\mathbf{K}\) -模同构

\[
\operatorname{Hom} _ {(\mathbf {A}, \mathbf {A})} (\mathfrak{J} , \mathrm{E}) \rightarrow \operatorname{Hom} _ {\mathbf {A} \otimes_ {\mathbf {K}} \mathbf {A}} (\mathfrak{J} , \mathrm{E}).
\]

另一方面， \(\mathfrak{J}\mathbf{E} = \{0\}\) ，因为元素 \(1 \otimes x - x \otimes 1\) 生成 \(\mathfrak{J}\) （作为 \((\mathbf{A} \otimes_{\mathbb{K}} \mathbf{A})\) -模；第10号，引理1），并且对所有 \(z \in \mathbf{E}\) ，

\[
(1 \otimes x - x \otimes 1) z = 0
\]

这是根据 E 上 \((\mathbf{A} \otimes_{\mathbf{K}} \mathbf{A})\) -模结构的定义。由于 \(\mathfrak{J}\) 包含在 \((\mathbf{A} \otimes_{\mathbf{K}} \mathbf{A})\) -模 E 的零化子中，且 \(((\mathbf{A} \otimes_{\mathbf{K}} \mathbf{A})/\mathfrak{J})\) -模 E 上的模结构按定义正是最初给定的 E 上的 A-模结构，因此，考虑到

\[
\mathfrak{J} \otimes_ {\mathbf {K}} ((\mathbf {A} \otimes_ {\mathbf {K}} \mathbf {A}) / \mathfrak{J})
\]

到 \(\mathfrak{J} / \mathfrak{J}^2\) （§ 4，第 1 号，命题 1 的推论 1）的典范同构，存在一个典范的 K-模同构

\[
\operatorname{Hom} _ {\mathbf {A} \otimes_ {\mathbf {K}} \mathbf {A}} (\mathfrak{J} , \mathrm{E}) \rightarrow \operatorname{Hom} _ {\mathbf {A}} (\mathfrak{J} / \mathfrak{J}^ {2}, \mathrm{E}).
\]

考虑到第10节命题17，可见我们已证明以下命题：

命题18。设 A 是一个交换 K-代数，且 \(\mathfrak{J}\) 是这样一个理想：它是满射典范同态 \(m: \mathbf{A} \otimes_{\mathbf{K}} \mathbf{A} \to \mathbf{A}\) 的核，于是 A 同构于 \((\mathbf{A} \otimes_{\mathbf{K}} \mathbf{A}) / \mathfrak{J}\) ，且 \(\mathfrak{J} / \mathfrak{J}^2\) 典范地具有一个 A-模结构。设 \(d_{\mathbf{A} / \mathbf{K}}: \mathbf{A} \to \mathfrak{J} / \mathfrak{J}^2\) 是使每个 \(x \in \mathbf{A}\) 对应 \(x \otimes 1 - 1 \otimes x\) 模 \(\mathfrak{J}^2\) 的同余类的 K-线性映射。映射 \(d_{\mathbf{A} / \mathbf{K}}\) 是一个 K-导子，并且对于每个 A-模 E 和每个 K-导子 D： \(\mathbf{A} \to \mathbf{E}\) ，存在且仅存在一个 A-线性映射 \(g: \mathfrak{J} / \mathfrak{J}^2 \to \mathbf{E}\) ，使得 \(\mathrm{D} = g \circ d_{\mathbf{A} / \mathbf{K}}\) .

A-模 \(\mathfrak{J} / \mathfrak{J}^2\) 称为 A 的 K-微分模，记作 \(\Omega_{\mathbf{K}}(\mathbf{A})\) ；对所有 \(x \in \mathbf{A}\) ， \(d_{\mathbf{A} / \mathbf{K}}(x)\) （也记作 \(dx\) ）称为 \(x\) 的微分；已经看到（第10号，引理1），元素 \(d_{\mathbf{A} / \mathbf{K}}(x)\) （ \(x \in \mathbf{A}\) ）构成 A-模 \(\Omega_{\mathbf{K}}(\mathbf{A})\) 的一个生成系。命题18表明映射 \(g \mapsto g \circ d_{\mathbf{A} / \mathbf{K}}\) 是一个典范 A-模同构

\[
\phi_ {\mathbf {A}}: \operatorname{Hom} _ {\mathbf {A}} (\Omega_ {\mathbf {K}} (\mathbf {A}), \mathrm{E}) \rightarrow \mathrm{D} _ {\mathbf {K}} (\mathrm{A}, \mathrm{E})
\]

（ \(D_{K}(A, E)\) 上的 A-模结构由第5号命题7定义）。

有序对 \((\Omega_{\mathbf{K}}(\mathbf{A}), d_{\mathbf{A}/\mathbf{K}})\) 因而是下述泛映射问题的解，其中 \(\Sigma\) 是 A-模结构的种类，而 \(\alpha\) -映射是从 A 到 A-模的 K-导子（集合论，IV，§3，第1号）。

例。设 M 是一个 K-模；由第9号命题14可知，对于每个 S(M)-模 E，映射 D ↦ D \textbar{} M 定义了 D\_K(S(M), E) 到 Hom\_K(M, E) 的一个 S(M)-模同构；另一方面，由于 E 是 S(M)-模，Hom\_K(M, E) 典范同构于

\[
\operatorname{Hom} _ {\mathbf {S} (\mathbf {M})} (\mathbf {M} \otimes_ {\mathbf {K}} \mathbf {S} (\mathbf {M}), \mathbf {E}),
\]

M 到 E 的每个 K-同态都能唯一地表示为 \(x \mapsto h(x \otimes 1)\) 的形式，其中

\[
h \in \operatorname{Hom} _ {\mathbf {S} (\mathbf {M})} (\mathbf {M} \otimes_ {\mathbf {K}} \mathbf {S} (\mathbf {M}), \mathbf {E})
\]

(II, § 5, 第 1 号)。设 \(\mathbf{D}_0\) 为 K-导子 \(\mathbb{S}(\mathbf{M}) \to \mathbf{M} \otimes_{\mathbb{K}} \mathbb{S}(\mathbf{M})\) ，它限制在 \(\mathbf{M}\) 上是典范同态 \(x \mapsto x \otimes 1\) ；因此每个 K-导子 \(\mathbf{D}: \mathbb{S}(\mathbf{M}) \to \mathbf{E}\) 都可以唯一地写成 \(h \circ \mathbf{D}_0\) ，其中

\[
h \in \operatorname{Hom} _ {\mathbf {S} (\mathbf {M})} (\mathbf {M} \otimes_ {\mathbf {K}} \mathbf {S} (\mathbf {M}), \mathbf {E}).
\]

由泛映射问题解的唯一性可知，存在唯一的S(M)-模同构

\[
\omega : \mathbf {M} \otimes_ {\mathbf {K}} \mathbf {S} (\mathbf {M}) \rightarrow \Omega_ {\mathbf {K}} (\mathbf {S} (\mathbf {M}))
\]

，使得 \(\mathbf{D}_0\circ \omega = d_{\mathbb{S}(\mathbf{M}) / \mathbb{K}}\) ；换言之，对所有 \(x\in \mathbf{M}\) \(\omega (x\otimes 1) = dx.\) 特别是，如果 \(\mathbf{M}\) 是自由 K-模，带有基 \((e_{\lambda})_{\lambda \in \mathbf{L}},\Omega_{\mathbb{K}}(\mathbb{S}(\mathbf{M}))\) 便是自由 S(M)-模，其基由微分 \(de_{\lambda}\) 的集合构成。特别考虑 \(\mathrm{L} = [1,n]\) 的情形，此时 \(\mathbb{S}(\mathbf{M})\) 与多项式代数 \(\mathbf{K}[\mathbf{X}_1,\dots ,\mathbf{X}_n]\) 等同（§ 6，第 6 号）；对每个多项式 \(\mathrm{P}\in \mathrm{K}[\mathrm{X}_1,\dots ,\mathrm{X}_n]\) ，我们可以唯一地写出

\[
d \mathrm{P} = \sum_ {i = 1} ^ {n} \mathrm{D} _ {i} \mathrm{P}. d \mathrm{X} _ {i}
\]

其中 \(D_{i}P \in K[X_{1}, \ldots, X_{n}]\) ，并且，根据上述内容，映射 \(P \mapsto D_{i}P\) 是 \(K[X_{1}, \ldots, X_{n}]\) 的 K-导子，对应于 \(\Omega_{K}(S(M))\) 关于基 \((dX_{i})\) 的坐标形式；我们也写作 \(\frac{\partial P}{\partial X_{i}}\) 以代替 \(D_{i}P\) ，并称之为 P 关于 \(X_{i}\) 的偏导数。

